\section{An exact integral for small Gram determinants}
\label{gv-section}

The geometric estimate for a degenerate orbital integral contains
the volume of tuples whose projected Gram determinant is small.
This volume admits a short proof by successive orthogonal projections.
We compute the relevant negative moments exactly, recover the submitted
estimate with a smaller constant, and also prove the optimal power
of the small determinant.  The passage from flag coordinates to this
Euclidean problem is a separate geometric step.
\cite{alpo-uniform-source}

Give $\mathbb C^n$ its usual Hermitian norm and real Lebesgue measure.
Let $P$ be an orthogonal projection of complex rank $d$, and let
$1\le j\le d\le n$.  For $v=(v_1,\ldots,v_j)$ put
\[
 \Gamma_P(v)=\Bigl[\det\bigl(\langle Pv_a,Pv_b\rangle\bigr)_{a,b=1}^j\Bigr]^{1/2}.
                                                               \tag{GV1}
\]
The determinant is real and nonnegative; its square root is the
nonnegative one.  Equivalently, $\Gamma_P$ is the norm of the
exterior product $Pv_1\wedge\cdots\wedge Pv_j$.  For $R>0$ let
\[
 \mathcal B_R=\{(v_1,\ldots,v_j):\|v_i\|\le R\text{ for every }i\}.
\]
We use the conventions $0!=1$ and that an empty product is one.

\subsection{One projected vector}

\begin{lemma}\label{gv-one-vector}
Let $Q$ be an orthogonal projection of rank $h$ on $\mathbb C^n$,
where $1\le h\le n$.  If $a$ is an integer with $0\le a<h$, then
\[
 \int_{\|v\|\le R}\|Qv\|^{-2a}\,dv
 =\frac{\pi^n R^{2(n-a)}}{(n-a)!}
       \frac{(h-a-1)!}{(h-1)!}.                               \tag{GV2}
\]
For $a\ge h$ the same integral is infinite.
\end{lemma}

\begin{proof}
The volume of a ball of radius $R$ in $\mathbb C^m$ is
$\pi^mR^{2m}/m!$.  One way to obtain the constant is to integrate
$e^{-\|x\|^2}$ first as a product of $m$ planar Gaussian integrals,
and then radially: the two values are $\pi^m$ and $m!$ times the
unit-ball volume.  Consequently radial integration in $\mathbb C^h$
uses the factor $2\pi^h r^{2h-1}/(h-1)!$.

After a unitary change of coordinates write $v=(u,w)$ with
$u\in\mathbb C^h$, $w\in\mathbb C^{n-h}$ and $Qv=(u,0)$.
Tonelli's theorem and the ball-volume formula give
\[
 \frac{2\pi^n}{(h-1)!(n-h)!}
   \int_0^R r^{2h-1-2a}(R^2-r^2)^{n-h}\,dr.
\]
This formula also covers $n=h$, since the zero-dimensional
ball has volume one.  Substituting $u=r^2/R^2$ changes the
last expression into
\[
 \frac{\pi^nR^{2(n-a)}}{(h-1)!(n-h)!}
   \int_0^1 u^{h-a-1}(1-u)^{n-h}\,du.
\]
For positive integers $b$ and nonnegative integers $c$,
successive integration by parts gives
\[
 \int_0^1 u^{b-1}(1-u)^c\,du
      =\frac{(b-1)!c!}{(b+c)!}.
\]
Use $b=h-a$, $c=n-h$ to obtain (GV2).
If $a\ge h$, restrict $w$ to a ball of radius $R/2$ and
$u$ to a sufficiently small ball.  The radial integral in
$u$ then contains a positive multiple of
$\int_0^{R/2}r^{2h-1-2a}\,dr$, which diverges.
\end{proof}

\subsection{Negative moments and a uniform sublevel estimate}

\begin{theorem}\label{gv-negative-moment}
If $a$ is an integer with $0\le a\le d-j$, then
\[
 \int_{\mathcal B_R}\Gamma_P(v)^{-2a}\,dv
 =\left(\frac{\pi^nR^{2(n-a)}}{(n-a)!}\right)^j
       \prod_{i=1}^j\frac{(d-i-a)!}{(d-i)!}.                  \tag{GV3}
\]
In particular the integral is finite.  The set $\Gamma_P=0$
has Lebesgue measure zero.  If $a\ge d-j+1$, the negative
moment in (GV3) is infinite.
\end{theorem}

\begin{proof}
Almost every projected tuple is linearly independent.  Indeed,
once $Pv_1,\ldots,Pv_{i-1}$ are independent, the condition
$Pv_i\in\operatorname{span}(Pv_1,\ldots,Pv_{i-1})$ describes a
proper complex linear subspace of $\mathbb C^n$, of codimension
$d-i+1\ge1$.  It has real Lebesgue measure zero.  Induction
and Fubini prove the assertion for the entire tuple.

For an independent partial tuple let $Q_i$ be the orthogonal
projection onto
\[
 \operatorname{im}P\cap
        \operatorname{span}(Pv_1,\ldots,Pv_{i-1})^\perp.
\]
Its rank is $d-i+1$.  Gram--Schmidt orthogonalization gives
\[
 \Gamma_P(v_1,\ldots,v_i)
   =\Gamma_P(v_1,\ldots,v_{i-1})\,\|Q_iv_i\|,
 \qquad \Gamma_P(\varnothing)=1.                              \tag{GV4}
\]
Integrate first in $v_j$.  Lemma~\ref{gv-one-vector} shows
that its contribution is
\[
 \frac{\pi^nR^{2(n-a)}}{(n-a)!}
       \frac{(d-j-a)!}{(d-j)!},
\]
independently of the preceding independent vectors.  Repeat
with $v_{j-1},\ldots,v_1$.  This proves (GV3).  All exchanges
are valid by Tonelli, even before finiteness is established;
the exceptional dependent tuples form a null set.  When $a=0$
the integrand is interpreted as one on that null set.

If $a\ge d-j+1$, then the conditional integral in $v_j$
is infinite for every independent preceding tuple, by the
last assertion of Lemma~\ref{gv-one-vector}.  Such preceding
tuples have full measure in a product of balls of positive
measure.  Tonelli therefore gives divergence.
\end{proof}

\begin{corollary}\label{gv-sublevel}
Let $1\le j<d\le n$, and write $s=d-j$.  Then for $t\ge0$,
\[
 \begin{split}
 &\operatorname{vol}\{v\in\mathcal B_R:\Gamma_P(v)\le t\}\\
 &\quad\le
  \frac{\pi^{nj}R^{2j(n-s)}t^{2s}}{((n-s)!)^j}
       \prod_{b=0}^{j-1}\frac{b!}{(s+b)!}.                    \end{split}
 \tag{GV5}
\]
The same upper bound holds if the rank of $P$ is at least $d$.
It also holds with every $v_i$ replaced by $Av_i$ both in the
norm conditions and in the Gram determinant, for any complex
linear automorphism $A$ with $|\det A|=1$.
\end{corollary}

\begin{proof}
For $t>0$, outside the null set $\Gamma_P=0$,
\[
        {\bf1}_{\{\Gamma_P\le t\}}
                    \le t^{2s}\Gamma_P^{-2s}.
\]
Integrating and applying (GV3) with $a=s$ proves (GV5).
For $t=0$ the null-set assertion proves the result directly.

If $\operatorname{rank}P\ge d$, choose a $d$-dimensional
subspace of its image and let $P_0$ be the orthogonal projection
onto it.  Then $\Gamma_{P_0}(v)\le\Gamma_P(v)$.  For example,
in coordinates adapted to the two images, Cauchy--Binet writes
the squared Gram determinant as a sum of squared absolute
values of minors; the smaller projection retains a subset
of these terms.  The sublevel set for $P$ is therefore
contained in the corresponding sublevel set for $P_0$.

Finally the real Jacobian of a complex linear map $A$ is
$|\det A|^2$.  The simultaneous substitution $w_i=Av_i$
preserves product Lebesgue measure when $|\det A|=1$.
\end{proof}

\subsection{The optimal power of the threshold}

The divergent moment at the next integer does not prevent
a stronger sublevel estimate.  To obtain it, integrate the
last perpendicular ball directly before applying a negative
moment to the preceding vectors.

\begin{theorem}\label{gv-optimal-power}
Let $1\le j\le d\le n$ and $q=d-j+1$.  Then for $t\ge0$,
\[
 \begin{split}
 &\operatorname{vol}\{v\in\mathcal B_R:\Gamma_P(v)\le t\}\\
 &\quad\le
  \frac{\pi^{nj}R^{2j(n-q)}t^{2q}}{q!\,((n-q)!)^j}
       \prod_{i=1}^{j-1}\frac{(j-i-1)!}{(d-i)!}.              \end{split}
 \tag{GV6}
\]
The extensions to rank at least $d$ and to $|\det A|=1$
in Corollary~\ref{gv-sublevel} remain valid.
For a projection of rank exactly $d$, no exponent larger
than $2q$ can replace $2q$ in an estimate with a finite
constant independent of $t\in(0,1)$.
\end{theorem}

\begin{proof}
First suppose $j\ge2$, and fix an independent preceding tuple
$v'=(v_1,\ldots,v_{j-1})$.  Write $\gamma=\Gamma_P(v')$.
The decomposition defined by $Q_j$ splits $v_j$ into a vector
in $\mathbb C^q$ and a vector in $\mathbb C^{n-q}$.
By (GV4) the first vector has norm at most $t/\gamma$.
The second has norm at most $R$.  Dropping their coupled
ball constraint gives the conditional volume bound
\[
 \frac{\pi^q}{q!}(t/\gamma)^{2q}
       \frac{\pi^{n-q}R^{2(n-q)}}{(n-q)!}.
\]
For the preceding $j-1$ vectors the exponent $a=q$ satisfies
$a=d-(j-1)$, so (GV3) applies.  Multiplication of its value
by the displayed factor proves (GV6).  For $j=1$, the same
two-ball estimate directly proves the formula, with the
product empty.  The case $t=0$ follows from the null-set
assertion.  The two extensions follow exactly as before.

To prove the optimality assertion, choose an orthonormal
basis of $\operatorname{im}P$.  When $j\ge2$, restrict
$v_1,\ldots,v_{j-1}$ to sufficiently small neighborhoods
of $(R/4)e_1,\ldots,(R/4)e_{j-1}$.  This gives a measurable
set of positive volume on which the projected vectors are
independent, all norms are less than $R/2$, and
\[
                  0<\gamma\le C=(R/2)^{j-1}.
\]
Decompose the last vector orthogonally into its components
in the preceding projected span, its $q$-dimensional
complement within $\operatorname{im}P$, and $\ker P$.
Restrict the first and third components to balls of radius
$R/4$, and the second to a ball of radius $t/C$.  For
$0<t\le CR/4$, all these last vectors have norm less
than $R$, and (GV4) gives $\Gamma_P\le t$.
The conditional volume is a fixed positive constant times
$t^{2q}$, independently of the chosen preceding tuple.
After integration over that tuple, the full sublevel
volume is at least $c\,t^{2q}$ for some $c>0$.
For $j=1$ use the same decomposition with the preceding
span absent and $C=1$.

If a bound $C' t^{2q+\delta}$ with $\delta>0$ held for
all sufficiently small $t$, division by $t^{2q}$ would
give $c\le C't^\delta$, a contradiction as $t\to0$.
\end{proof}

\subsection{The two volume bounds in the orbital calculation}

For direct comparison with the submitted estimates, put
$n=d+c$, $s=d-j\ge1$, and $R=M\ge1$.  Each factorial
ratio in (GV5) satisfies
\[
             \frac{b!}{(s+b)!}\le\frac1{s!}.
\]
Also $s!\ge(s/e)^s$: the term $s^s/s!$ of the exponential
series is at most $e^s$.  Hence (GV5) implies
\[
 \begin{split}
 &\operatorname{vol}\{v:\|Av_i\|\le M,\quad\Gamma_P(Av_1,\ldots,Av_j)\le t\}\\
 &\quad\le
   \frac{\pi^{j(j+c)}M^{2j(j+c)}}{((j+c)!)^j}
        \left(\frac{\pi e}{s}\right)^{js}t^{2s},
                                                               \end{split}
 \tag{GV7}
\]
provided $\operatorname{rank}P\ge d$ and $|\det A|=1$.

For the degenerate orbital integral, take
\[
 n=k+1,\qquad d=k-1,\qquad c=2,\qquad
 s=k-1-j,\qquad t=2y,
\]
where $1\le j$ and $j+2\le k$.  The source's middle
isometry $m$ satisfies $m^*H m=H$ for a nondegenerate
Hermitian matrix $H$.  Taking determinants gives
$|\det m|^2=1$, so $A=m^{-1}$ has the required Jacobian.
The sublevel volume is therefore at most
\[
 \frac{\pi^{j^2+2j}M^{2j^2+4j}}{((j+2)!)^j}
     \left(\frac{\pi e}{k-1-j}\right)^{j(k-1-j)}
     (2y)^{2(k-1-j)}.                                        \tag{GV8}
\]
The submitted bound has the same expression without the
factorial denominator and with an additional factor $2$.
Thus (GV8) improves its constant by at least
$2((j+2)!)^j$, with exactly the same hypotheses and powers
of $M,y$.  It also proves that the zero Gram locus is null.

The other orbital volume bound uses
$d=k-2$, $c=3$, $j=r$, and $s=k-2-r\ge1$.
The same argument gives
\[
 \frac{\pi^{r^2+3r}M^{2r^2+6r}}{((r+3)!)^r}
      \left(\frac{\pi e}{k-2-r}\right)^{r(k-2-r)}
      (2y)^{2(k-2-r)}.                                       \tag{GV9}
\]
This bounds the entire sublevel set and hence also its shell
$y<\Gamma_P\le2y$.  The submitted expression has a larger
factor $4M^{2r}$ and omits the factorial denominator.
Its assumed range $2r+6\le k$ is stronger than needed here.

The sharper exponent in (GV6) is also available for both
applications.  Using it in the complete orbital bound would
change the subsequent determinant powers and their comparison
with the lattice count, which must be tracked separately.
For the present reconstruction, (GV8)--(GV9) already supply
the exact volume estimates required by the existing argument.
They do not establish the flag-coordinate identity, the
coordinate-size bounds, or the normalization of the common
unipotent stabilizer; those remain separate inputs to the
orbital calculation.


